% cSpell:locale en,pt-BR - Adicionado para suporte ao corretor ortográfico
% Extensão Code Spell Checker no VS Code (muito útil!)

% =====================================================================
%  MANUAL DO MODELO ABNT DO ICT/UNIFEI
%
%  Este arquivo NÃO é um ponto de partida. É o manual — o documento que
%  explica como usar o modelo, e que se lê compilado, em manual.pdf.
%
%  Para começar um trabalho, copie o arquivo de partida do seu tipo:
%  modelo-generico.tex, modelo-estagio.tex, modelo-tcc1.tex,
%  modelo-tcc2.tex, modelo-dissertacao.tex ou modelo-tese.tex. O capítulo 1
%  explica qual escolher e o que cada um pressupõe.
%
%  O manual também é exemplo de si mesmo: ele se organiza em Capitulos/,
%  uma pasta por capítulo, e demonstra \subimport em uso. Os arquivos de
%  partida demonstram a alternativa, tudo em um arquivo só.
%
%  Compile com:  latexmk manual.tex
% =====================================================================
\documentclass[12pt,a4paper,oneside]{book}

% O manual compila como generico porque precisa de algum tipo, e esse é o
% que menos promete sobre si mesmo. Os seis tipos, e o que muda entre eles,
% são assunto do capítulo 3 — não deste comentário.
\usepackage[tipo=generico]{UnifeiICTReport}
%\usepackage[utf8]{inputenc} % Removido pois XeLaTeX usa UTF-8 nativo
\usepackage{csquotes}
\usepackage{fontspec}
\usepackage{import}
\usepackage{lipsum} % Para gerar texto placeholder
\usepackage[style=abnt]{biblatex}
\usepackage{fancyvrb} % Para formatação avançada de verbatim
\usepackage{subcaption}
\usepackage{booktabs}

% \usepackage{showframe} % para debug de margens e layout
% \usepackage{lua-visual-debug} % para debug visual com LuaLaTeX

% Trechos em monoespaçado não hifenizam, e um nome de comando longo no fim da
% linha estoura a margem. O emergencystretch só é usado quando não há outra
% saída, então as demais linhas continuam justificadas como antes.
\setlength{\emergencystretch}{3em}

\makeglossaries % necessário para glossários e listas de abreviaturas
\subimport{Preambulo/}{lista-abreviaturas.tex}
\subimport{Preambulo/}{lista-simbolos.tex}

\bibliography{referencias.bib}

%--------------------------------------
% Informações de dados para CAPA e FOLHA DE ROSTO
%--------------------------------------

% Título, subtítulo e autores (compatibilidade com unifeiICT.sty)
\title{Modelo ABNT do ICT/Unifei}

% Comente ou remova a linha abaixo se não houver subtítulo.
% Escreva SÓ o texto do subtítulo, sem dois-pontos: a NBR 14724:2024 §4.1.1(d)
% exige que ele venha "precedido de dois-pontos, evidenciando a sua
% subordinação ao título", e o modelo insere essa pontuação sozinho na capa,
% na folha de rosto e na folha de aprovação.
\subtitle{manual de uso}

\author{Wendell F. S. Diniz}

% O manual não tem professor responsável no sentido usual — ele não é entregue a
% ninguém. A linha fica declarada com o nome de quem mantém o modelo, pelo mesmo
% motivo que \course permanece em Engenharia de Computação: registro de autoria.
%
% Deixá-la em branco não é opção enquanto a change fix-supervisor-guard-approval-sheet
% não entrar: \folhaaprovacao não consulta \if@supervisorpresent, ao contrário da
% folha de rosto, e emitiria o rótulo "Professor da disciplina" sem nome acima.
\supervisor{Prof. Dr. Wendell F. S. Diniz}

% \cosupervisor{Prof. Dra. Nome da Coorientadora}

% Em generico, \subject é a disciplina que o trabalho atende. O manual não
% atende disciplina alguma; o campo diz o que ele é.
\subject{Documentação do pacote UnifeiICTReport}

% Metadados adicionais usados pela folha de aprovação (\folhaaprovacao, mais
% abaixo). Data de aprovação e nota são opcionais; sem elas, a folha sai sem
% essas linhas, em vez de com campos em branco estranhos.
% \aprovacaodata{15 de dezembro de 2025}
% \notaaprovacao{9,5}

% Só para os tipos de monografia (tcc2, dissertacao, tese).
% A área de concentração é EXIGIDA em dissertacao e tese — sem ela a
% compilação para, nomeando o que falta. No tcc2 ela é opcional e rara: só
% quando o trabalho integra um projeto maior. A linha de pesquisa é sempre
% opcional.
% \areaconcentracao{Sistemas de Computação}
% \linhapesquisa{Engenharia de Software}

% Membros da banca examinadora — usados por tcc1 (quando há defesa), tcc2,
% dissertacao e tese. Declare apenas os membros EXTERNOS: o orientador já vem
% de \supervisor e entra na banca sozinho, sem precisar ser repetido aqui.
% Campos em branco ainda rendem um bloco corretamente formatado, para a folha
% preparada antes de os nomes serem conhecidos.
% No tcc1 a defesa é facultativa: sem nenhum \bancamembro declarado, a folha
% de aprovação simplesmente não é emitida.
% \bancamembro{Prof. Dr. Fulano de Tal}{Doutor em Engenharia}{Unifei}
% \bancamembro{Prof. Dra. Beltrana de Tal}{Doutora em Engenharia}{Unifei}


% Resumo na língua principal
\abstract{%
    Este manual descreve o pacote \LaTeX{} UnifeiICTReport, que formata trabalhos acadêmicos do
    Instituto de Ciências Tecnológicas da Universidade Federal de Itajubá conforme as normas da
    ABNT. Apresenta os seis tipos de documento que o pacote atende --- relatório de disciplina,
    relatório de estágio, projeto de pesquisa, monografia de conclusão de curso, dissertação e tese
    ---, a norma que rege cada um e o que o tipo declarado governa nos elementos pré-textuais.
    Documenta os comandos e ambientes disponíveis ao autor: metadados de capa e folha de rosto,
    folha de aprovação, resumos, divisões do texto, referências cruzadas, citações, ilustrações,
    tabelas, equações e elementos pós-textuais. Registra também o que o pacote não oferece e os
    pontos em que ele se afasta deliberadamente da norma. Destina-se a quem vai escrever um
    trabalho com o modelo, e não a quem mantém o pacote.
}
% Palavras-chave do resumo principal
\keywords{ABNT; \LaTeX{}; trabalhos acadêmicos; Unifei; ICT.}

% Resumo na língua secundária, geralmente inglês.
\abstractseclang{%
    This manual describes the UnifeiICTReport \LaTeX{} package, which formats academic work from
    the Institute of Technological Sciences of the Federal University of Itajubá according to ABNT
    standards. It presents the six document types the package supports --- course report,
    internship report, research project, undergraduate thesis, master's dissertation and doctoral
    thesis ---, the standard governing each one, and what the declared type controls in the
    pre-textual elements. It documents the commands and environments available to the author:
    title page metadata, approval sheet, abstracts, text divisions, cross-references, citations,
    illustrations, tables, equations and post-textual elements. It also records what the package
    does not provide and where it deliberately departs from the standard. It is intended for those
    writing a document with the template, not for those maintaining the package.
}
% Palavras-chave do resumo secundário
\keywordsseclang{ABNT; \LaTeX{}; academic work; Unifei; ICT.}

% ---
% Configurações de pacotes
% ---

\begin{document}

% Preâmbulo do documento (elementos pré-textuais)
\frontmatter

% ---
% Impressão da capa e folha de rosto
% ---

\maketitle

% ---
% Folha de aprovação (NBR 14724:2024 §4.2.1.3) — sem argumento: a forma dela
% vem do tipo declarado lá em cima, no \usepackage.
% No tcc1 ela só sai se houver banca declarada; nos demais tipos é obrigatória.
% ---
\folhaaprovacao

% ---
% Resumos
% ---
\makeabstracts

%----
%  Listas de figuras, tabelas, abreviaturas e símbolos
%  -- Estes elementos são opcionais e podem ser incluídos conforme a necessidade do relatório.
%  -- Para inclui-los, descomente as linhas abaixo. Para suprimir, deixe comentado.
% ----

\listoffigures

\listoftables

% Quadro e Gráfico (NBR 14724:2024 §5.8) têm lista própria, como Figura e
% Tabela. Descomente a que fizer sentido para o relatório; nenhum tipo usado
% no texto gera página em branco por a lista ficar comentada.
% \listofquadros
%
% \listofgraficos

\printglossary[type=\acronymtype,title={\acronymlistname}]

\printglossary[type=symbols,title={\symbolslistname}]

% ---
% Sumário
% ---
\tableofcontents

\mainmatter

% Os capítulos ficam em Capitulos/, uma pasta por capítulo. É esta a
% organização que o manual recomenda para trabalhos longos, e ele a segue para
% servir de exemplo dela. Comente uma linha para pular aquele capítulo durante
% a edição; descomente todas antes da compilação final.
\subimport{Capitulos/cap1/}{cap1.tex}
\subimport{Capitulos/cap2/}{cap2.tex}
\subimport{Capitulos/cap3/}{cap3.tex}
\subimport{Capitulos/cap4/}{cap4.tex}
\subimport{Capitulos/cap5/}{cap5.tex}
\subimport{Capitulos/cap6/}{cap6.tex}
\subimport{Capitulos/cap7/}{cap7.tex}

% ---
% Referências bibliográficas
% --- Imprime a bibliografia
% ---
\printbibliography

% ---
% Elementos pós-textuais são iniciados após \backmatter
% --- Compreende apêndices, anexos, glossários e índices
% --- São todos opcionais
% ---

\backmatter

% Apêndices: produzidos pelo próprio autor, para complementar a argumentação
% (NBR 14724:2024 §3.4). \apendices abre o grupo (uma entrada coletiva no
% sumário); cada \apendice{<rótulo>}{<título>} abre um apêndice, lettered A, B, ...
\apendices

\apendice{ape:questionario}{Questionário aplicado na pesquisa}
Este apêndice é de autoria do próprio autor do relatório — daí ser um
\emph{apêndice}, e não um anexo. Ilustrações dentro de um apêndice continuam a
mesma sequência numérica do restante do documento e aparecem normalmente nas
listas, como a \reffig{fig:questionario} a seguir demonstra.

\begin{figuraabnt}{fig:questionario}{Estrutura do questionário aplicado}
    \includegraphics[width=0.6\textwidth]{example-image}
    \fonte{o próprio autor}
\end{figuraabnt}

% Anexos: documentos de terceiros, incorporados como fundamentação, prova ou
% ilustração (§3.3) -- não escritos pelo autor do relatório.
\anexos

\anexo{ane:certificado}{Certificado de participação no evento}
Este anexo reproduz um documento de terceiros — o certificado emitido pela
organização do evento — anexado ao relatório como prova, e não redigido pelo
autor. Essa autoria é exatamente o que separa um \emph{anexo} de um
\emph{apêndice}: veja o \refcomp{}{ape:questionario}, de autoria própria, em
contraste com este.

\end{document}